\section{OSPF Area Healing Across ABRs}

\subsection{Objective}

The objective of this experiment was to validate OSPF area healing across ABRs.

The experiment tested whether traffic between two different OSPF areas could recover when the direct ABR-to-ABR/core link failed.

The traffic tested was:

\begin{verbatim}
hpe-h1 -> hpe-h2
10.0.61.2 -> 10.0.82.2
\end{verbatim}

In the topology, \texttt{hpe-h1} is behind Area 10, and \texttt{hpe-h2} is behind Area 20. Router \texttt{hpe-r3} acts as the ABR between Area 10 and Area 0, while \texttt{hpe-r4} acts as the ABR between Area 20 and Area 0.

\subsection{Link Failed}

The failed link was the direct core/ABR link between \texttt{hpe-r3} and \texttt{hpe-r4}.

The failed interface was:

\begin{verbatim}
hpe-r3 eth0
\end{verbatim}

Before the failure, \texttt{hpe-r3} reached the Area 20 host network through \texttt{hpe-r4}.

\begin{verbatim}
10.0.82.0/24
via 10.0.34.3 on eth0
\end{verbatim}

The old direct next-hop was:

\begin{verbatim}
10.0.34.3
\end{verbatim}

\subsection{Baseline Route Before Failure}

Before the failure, the route from \texttt{hpe-r3} toward \texttt{10.0.82.0/24} was:

\begin{verbatim}
10.0.82.0/24 unicast [ospf1] IA
via 10.0.34.3 on eth0
\end{verbatim}

The Linux kernel forwarding decision also used the same next-hop:

\begin{verbatim}
10.0.82.2 via 10.0.34.3 dev eth0
\end{verbatim}

This proves that before failure, Area 10 to Area 20 traffic used the direct \texttt{hpe-r3 -> hpe-r4} path.

\subsection{Healing Behaviour After Failure}

After \texttt{hpe-r3 eth0} was brought down, OSPF recalculated an alternate path.

The route-get sampler first showed the new forwarding path as:

\begin{verbatim}
10.0.82.2 via 10.0.23.2 dev eth3
\end{verbatim}

The Linux kernel route later showed two alternate next-hops:

\begin{verbatim}
10.0.82.0/24 proto bird metric 32
nexthop via 10.0.13.2 dev eth2 weight 1
nexthop via 10.0.23.2 dev eth3 weight 1
\end{verbatim}

The BIRD route table also showed the healed route:

\begin{verbatim}
10.0.82.0/24 unicast [ospf1] IA
via 10.0.13.2 on eth2 weight 1
via 10.0.23.2 on eth3 weight 1
\end{verbatim}

This proves that after the direct ABR link failed, traffic healed through alternate backbone paths via \texttt{hpe-r1} and \texttt{hpe-r2}.

\subsection{Measurement Results}

\begin{table}[h]
\centering
\begin{tabular}{|l|r|}
\hline
\textbf{Measurement} & \textbf{Result} \\
\hline
Route-get switch time & 80 ms \\
\hline
Kernel route switch time & 931 ms \\
\hline
BIRD route switch time & 962 ms \\
\hline
Ping packets transmitted & 150 \\
\hline
Ping packets received & 136 \\
\hline
Ping packets lost & 14 \\
\hline
Packet loss & 9.33333\% \\
\hline
\end{tabular}
\caption{OSPF area healing measurements}
\end{table}

The missing ICMP sequence numbers were:

\begin{verbatim}
11 12 13 14 15 16 17 18 19 20 21 22 23 24
\end{verbatim}

The ping interval used during the test was 0.1 seconds. This means the 14 lost packets represent a short disruption window of around 1.4 seconds during OSPF area healing.

\subsection{Final Restored Health}

After the failed link was restored, connectivity was checked again.

\begin{verbatim}
5 packets transmitted
5 received
0\% packet loss
\end{verbatim}

This confirms that the network returned to a healthy state after the link was restored.

\subsection{Interpretation}

The experiment shows three important things.

First, the direct ABR path between \texttt{hpe-r3} and \texttt{hpe-r4} was being used before the failure.

Second, when that direct link failed, OSPF did not leave Area 10 to Area 20 traffic broken. It recalculated alternate Area 0 backbone paths through \texttt{hpe-r1} and \texttt{hpe-r2}.

Third, after the link was restored, normal connectivity returned with 0\% packet loss.

The route-get result changed quickly, within 80 ms. The full kernel and BIRD route-table convergence took around 1 second. During this convergence window, 14 out of 150 high-frequency ping packets were lost.

\subsection{Conclusion}

The OSPF area healing experiment successfully demonstrated inter-area recovery across ABRs.

Before failure, traffic from Area 10 toward Area 20 used the direct path from \texttt{hpe-r3} to \texttt{hpe-r4} through next-hop \texttt{10.0.34.3}. After the \texttt{hpe-r3 eth0} link failed, OSPF recalculated the path and installed alternate next-hops through \texttt{10.0.13.2} and \texttt{10.0.23.2}.

The route-get sampler detected a forwarding change in 80 ms, while the full kernel and BIRD route-table updates completed in around 1 second. During the transition, 14 packets were lost in a high-frequency ping test, but after restoration the final health check showed 0\% packet loss.

This confirms that OSPF area healing worked correctly across ABRs in the lab topology.
